\def\tituloconcepto{Tomografía de estados}
% Si el capitulo se compila solo (sin main.tex), esta cabecera abre el documento.
% Antes de % Si el capitulo se compila solo (sin main.tex), esta cabecera abre el documento.
% Antes de % Si el capitulo se compila solo (sin main.tex), esta cabecera abre el documento.
% Antes de \input{cabecera} cada capitulo define \tituloconcepto.
\ifdefined\enmain\else
  \documentclass[11pt,a4paper]{article}
  \input{preambulo}
  \begin{document}
  \title{\tituloconcepto}
  \author{\autor}
  \date{\today}
  \maketitle
\fi
 cada capitulo define \tituloconcepto.
\ifdefined\enmain\else
  \documentclass[11pt,a4paper]{article}
  % Preambulo compartido por todos los documentos. Compilar con XeLaTeX/LuaLaTeX
% (o Tectonic). Si se usa Overleaf, elegir el compilador XeLaTeX.
\usepackage[spanish,es-noshorthands]{babel}
\usepackage{amsmath,amssymb}
\usepackage{booktabs,array}
\usepackage[a4paper,margin=2.5cm]{geometry}
\usepackage{xcolor}
\usepackage{listings}
\usepackage[most]{tcolorbox}
\usepackage{quantikz}
\usepackage[hidelinks]{hyperref}

% ---- Autor (editar a gusto) -------------------------------------------------
\newcommand{\autor}{Genesis Silesky}

% ---- Notacion de Dirac y atajos ---------------------------------------------
% quantikz ya define \ket y \bra; DeclareRobustCommand los redefine sin error.
\DeclareRobustCommand{\ket}[1]{\left|#1\right\rangle}
\DeclareRobustCommand{\bra}[1]{\left\langle#1\right|}
\DeclareRobustCommand{\braket}[2]{\left\langle#1\middle|#2\right\rangle}
\newcommand{\C}{\mathbb{C}}
\newcommand{\R}{\mathbb{R}}
\DeclareMathOperator{\Tr}{Tr}

% ---- Cajas de apoyo --------------------------------------------------------
% "concepto": definiciones previas que hacen falta para entender el capitulo.
\newtcolorbox{concepto}[1]{breakable,colback=blue!4,colframe=blue!55!black,
  fonttitle=\bfseries,title={Concepto previo: #1}}
% "nota": aclaraciones y advertencias.
\newtcolorbox{nota}[1]{breakable,colback=orange!6,colframe=orange!70!black,
  fonttitle=\bfseries,title={#1}}

% ---- Estilos de listados ----------------------------------------------------
\lstdefinestyle{py}{language=Python,basicstyle=\ttfamily\footnotesize,
  breaklines=true,postbreak=\mbox{\textcolor{gray}{$\hookrightarrow$}\space},
  numbers=left,numberstyle=\tiny\color{gray},frame=single,
  commentstyle=\color{green!40!black}\itshape,
  keywordstyle=\color{blue!70!black}\bfseries,
  stringstyle=\color{red!60!black},showstringspaces=false,tabsize=4,
  columns=fullflexible,keepspaces=true}
\lstdefinestyle{salida}{basicstyle=\ttfamily\footnotesize,breaklines=true,
  frame=leftline,backgroundcolor=\color{gray!8},columns=fullflexible,
  keepspaces=true}

  \begin{document}
  \title{\tituloconcepto}
  \author{\autor}
  \date{\today}
  \maketitle
\fi
 cada capitulo define \tituloconcepto.
\ifdefined\enmain\else
  \documentclass[11pt,a4paper]{article}
  % Preambulo compartido por todos los documentos. Compilar con XeLaTeX/LuaLaTeX
% (o Tectonic). Si se usa Overleaf, elegir el compilador XeLaTeX.
\usepackage[spanish,es-noshorthands]{babel}
\usepackage{amsmath,amssymb}
\usepackage{booktabs,array}
\usepackage[a4paper,margin=2.5cm]{geometry}
\usepackage{xcolor}
\usepackage{listings}
\usepackage[most]{tcolorbox}
\usepackage{quantikz}
\usepackage[hidelinks]{hyperref}

% ---- Autor (editar a gusto) -------------------------------------------------
\newcommand{\autor}{Genesis Silesky}

% ---- Notacion de Dirac y atajos ---------------------------------------------
% quantikz ya define \ket y \bra; DeclareRobustCommand los redefine sin error.
\DeclareRobustCommand{\ket}[1]{\left|#1\right\rangle}
\DeclareRobustCommand{\bra}[1]{\left\langle#1\right|}
\DeclareRobustCommand{\braket}[2]{\left\langle#1\middle|#2\right\rangle}
\newcommand{\C}{\mathbb{C}}
\newcommand{\R}{\mathbb{R}}
\DeclareMathOperator{\Tr}{Tr}

% ---- Cajas de apoyo --------------------------------------------------------
% "concepto": definiciones previas que hacen falta para entender el capitulo.
\newtcolorbox{concepto}[1]{breakable,colback=blue!4,colframe=blue!55!black,
  fonttitle=\bfseries,title={Concepto previo: #1}}
% "nota": aclaraciones y advertencias.
\newtcolorbox{nota}[1]{breakable,colback=orange!6,colframe=orange!70!black,
  fonttitle=\bfseries,title={#1}}

% ---- Estilos de listados ----------------------------------------------------
\lstdefinestyle{py}{language=Python,basicstyle=\ttfamily\footnotesize,
  breaklines=true,postbreak=\mbox{\textcolor{gray}{$\hookrightarrow$}\space},
  numbers=left,numberstyle=\tiny\color{gray},frame=single,
  commentstyle=\color{green!40!black}\itshape,
  keywordstyle=\color{blue!70!black}\bfseries,
  stringstyle=\color{red!60!black},showstringspaces=false,tabsize=4,
  columns=fullflexible,keepspaces=true}
\lstdefinestyle{salida}{basicstyle=\ttfamily\footnotesize,breaklines=true,
  frame=leftline,backgroundcolor=\color{gray!8},columns=fullflexible,
  keepspaces=true}

  \begin{document}
  \title{\tituloconcepto}
  \author{\autor}
  \date{\today}
  \maketitle
\fi


\section{Tomografía de estados: reconstruir un estado a partir de mediciones}

\subsection{Objetivo}
Entender cómo se lee por completo el estado de uno o dos qubits, midiendo en
varios ejes y reconstruyendo la matriz densidad $\rho$.

\subsection{Conceptos previos}
\begin{concepto}{esfera de Bloch, matrices de Pauli y entrelazamiento}
Ver capítulo~10 (esfera de Bloch), capítulo~04 (medición y shots) y el
capítulo~12 (producto tensorial y entrelazamiento). Una \emph{matriz densidad}
$\rho$ es otra forma de describir un estado; para un estado puro es
$\ket\psi\bra\psi$.
\end{concepto}

\subsection{Fundamento teórico}
Al terminar un cálculo hay que leer el estado de los qubits. La \textbf{tomografía
cuántica} es la técnica para leerlo completo.

\paragraph{Un qubit.} Las coordenadas del qubit en la esfera de Bloch son
$(a,b,c)=(\langle\sigma_x\rangle,\langle\sigma_y\rangle,\langle\sigma_z\rangle)$,
los promedios de los momentos angulares de espín en cada eje. Medir colapsa el
estado y solo da el resultado de un autoestado, pero la información queda en las
probabilidades:
\begin{equation}
  P(0)=\frac{1+c}{2},\qquad P(1)=\frac{1-c}{2}\ \Longrightarrow\ c=P(0)-P(1).
\end{equation}
Repitiendo muchas veces se estima $c$. Con $(a,b,c)$ se reconstruye
\begin{equation}
  \rho=\frac12\big(I+a\,\sigma_x+b\,\sigma_y+c\,\sigma_z\big).
\end{equation}

\paragraph{Medir $x$ e $y$.} Algunos sistemas solo pueden medir en el eje $z$.
Para medir $\langle\sigma_x\rangle$ y $\langle\sigma_y\rangle$ se aplican antes
unos pulsos de observación $\pi/2$ que giran el eje que se quiere medir hasta
$z$. El cuadro completo se obtiene solo al reunir los tres ejes
perpendiculares (una ``vista en tres ángulos'').

\paragraph{Varios qubits.} Para $n$ qubits, $\rho$ se expande en los productos
de matrices de Pauli ($\sigma_0=I$, $\sigma_1=\sigma_x$, $\sigma_2=\sigma_y$,
$\sigma_3=\sigma_z$). Para dos qubits:
\begin{equation}
  \rho=\frac14\sum_{i,j=0}^{3}c_{ij}\,\sigma_i\otimes\sigma_j,\qquad c_{00}=1 ,
\end{equation}
una suma de 16 términos, de los cuales $c_{00}=1$ es la normalización y hay
que medir los otros 15. Los términos como $\sigma_x\otimes I$ son el momento
angular de un solo qubit; los términos como $\sigma_z\otimes\sigma_z$ son la
\emph{correlación} entre los dos. Para más qubits el principio es el mismo y solo
aumenta el número de mediciones.

\begin{nota}{Por qué no basta medir cada qubit por separado}
Si hubiera entrelazamiento, el estado completo no se puede obtener midiendo
cada qubit por separado y juntando los resultados, porque parte de la información
está en las correlaciones.
\end{nota}

\begin{nota}{Diferencia con el sistema de resonancia magnética}
En un sistema de resonancia magnética nuclear (RMN) las mediciones son sobre un
\emph{conjunto} de moléculas y la señal medida (el \emph{FID}) trae a la vez
información de dos ejes, por lo que bastan 6 observaciones para dos qubits. En
este capítulo se simula un computador que mide solo en $z$: el principio
(reconstruir $\rho$ con promedios de Pauli) es el mismo, pero se hacen más
mediciones (el código mide cada uno de los 15 coeficientes por separado).
\end{nota}

\subsection{Código (completamente comentado)}
\lstinputlisting[style=py]{../../fuentes/13_tomografia.py}

\subsection{Explicación celda por celda}
\begin{description}
  \item[Celda 1] Matrices de Pauli y un simulador con semilla fija.
  \item[Celda 3] \texttt{promedio} estima $\langle\sigma_{p_0}\otimes\sigma_{p_1}\otimes\cdots\rangle$:
    añade los giros que llevan cada eje a $z$ (\texttt{h} para $x$;
    \texttt{sdg} y \texttt{h} para $y$), mide y calcula el promedio con signo
    $+1$ o $-1$ según haya un número par o impar de unos.
  \item[Celda 4] Reconstruye $\rho$ de un qubit con $100$, $1000$ y $10\,000$
    mediciones y calcula la fidelidad con el estado real.
  \item[Celda 5] \texttt{tomografia\_2\_qubits} mide los 15 coeficientes
    distintos de $c_{00}$ y arma $\rho$.
  \item[Celda 6] Aplica la tomografía a un estado de Bell.
  \item[Celda 7] Compara con el producto $\frac I2\otimes\frac I2$ de las
    tomografías por separado.
\end{description}

\subsection{Salida obtenida}
\lstinputlisting[style=salida]{../../fuentes/salidas/13_tomografia.txt}
Se usa semilla fija, así que los números se repiten; con otra semilla los
valores de las mediciones cambian ligeramente.

\subsection{Interpretación}
Con $100$ mediciones el módulo de $(a,b,c)$ sale $1.044$, mayor que $1$: el
ruido estadístico produce una $\rho$ que no es un estado válido, y por eso la
``fidelidad'' pasa de $1$. Con $10\,000$ mediciones el módulo es $0.993$ y la
fidelidad $0.9966$. Para el estado de Bell los únicos coeficientes distintos de
cero son $c_{xx}=1$, $c_{yy}=-1$ y $c_{zz}=1$ (más $c_{00}=1$), y la fidelidad
con el Bell ideal es $1$. En cambio, el producto de las tomografías de cada
qubit por separado da $I/4$, con fidelidad de solo $0.25$: las correlaciones
eran justo lo que faltaba.

\subsection{Ejercicios}
\begin{enumerate}
  \item Haz la tomografía del estado $\ket+$ y verifica que
    $(a,b,c)\approx(1,0,0)$.
  \item Haz la tomografía de dos qubits en el estado $\ket{00}$.
    ¿Qué coeficientes son distintos de cero?
  \item Cuenta cuántos coeficientes habría que medir para 3 qubits.
\end{enumerate}

\subsection{Errores comunes}
\begin{itemize}
  \item Esperar un resultado exacto con pocas mediciones: el error baja
    como $1/\sqrt N$ (capítulo~04).
  \item Confundir el orden de las letras: en \texttt{promedio("XZ", ...)} la
    primera letra es el qubit 0 y la segunda el qubit 1.
\end{itemize}

\ifdefined\enmain\else
  \end{document}
\fi

